% !TEX root = ../main.tex
\section{Discussion}
\label{sec:discussion}

We document three findings: agentic-ready tools attract far more adoption
(up to 72\% within the same developer), this association appears on both the
extensive and conditional intensive margins, and we find no robust evidence
that adoption is more concentrated among agentic-ready tools. Together they
describe the adoption and concentration patterns associated with the platform's
readiness boundary for machine consumers.

%% ─────────────────────────────────────────────
\subsection{The Adoption Payoff of Agentic Readiness}
\label{sec:discussion:gap}

Within the same developer, agentic-ready tools attract a large and consistently
positive adoption gap---40\% to 72\% on the log-plus-one scale depending on
weighting (72\% under our preferred equal-developer-weighted estimator), and
several-fold larger on the count scale. This is the paper's central finding:
within a developer, the tools configured for agent-initiated use attract
substantially more adoption
than those that are not. The identification strategy isolates readiness, not realized agent activity.

Two pathways can produce this association, and our cross-sectional data cannot
separate them (\autoref{fig:mechanism}). Under an \emph{agent-mediated}
account, agentic-ready tools enter the feasible set agents can discover and
invoke, and this routing raises their adoption. Under a
\emph{developer-mediated} account, the architectural choices that make a tool
agentic-ready may coincide with developer investment or appeal to human
adopters independent of agents, so readiness is associated with adoption even
without any agent-side effect. Both accounts are consistent with the
extensive-margin finding.

Agentic readiness is the \emph{architectural} component of the platform's
agent-access admission, not a proxy we chose. The platform defines limited
permissions and non-standby as necessary architectural conditions for
agentic-payment eligibility; the \texttt{allowsAgenticUsers} filter and the MCP
\texttt{search-actors} tool surface agent-eligible tools to agent discovery, and
agent invocation of ineligible tools is rejected (\autoref{sec:context}). The
treatment therefore measures architectural admission to the agent-accessible
set---a necessary condition for agent use, distinct from the developer-level KYC
and approval-level curation gates the platform applies on top. What remains open
is the \emph{mechanism}: whether the adoption gap reflects agent routing,
correlated tool-level selection---task type, maturity, target audience---or both.

The estimate nonetheless survives excluding whale developers and
propensity-score weighting, so the association is not a small-developer
artifact. The remaining interpretive boundaries---unobserved agent use,
within-developer selection, the identifying subsample, and the temporal
mismatch---are detailed in \autoref{sec:discussion:limitations}.

%% ─────────────────────────────────────────────
\subsection{Within-Group Adoption Inequality}
\label{sec:discussion:concentration}

The paper's second finding is an asymmetry: agentic readiness is associated
with \emph{which tools get adopted}, but not with a \emph{detectable increase
in concentration}. We find no robust evidence that adoption is more
concentrated among agentic-ready tools: Gini coefficients are nearly identical
(0.935 versus 0.946), within categories agentic-ready tools are more
concentrated in only 7 of 15, and a sample-size-equalized comparison shows the
upper-tail gap (top-1\% 76.8\% versus 74.3\%) is too imprecisely estimated to
be informative. These estimates compare concentration within observed readiness sets; market-level counterfactual effects require a different design.

For platform operators, this asymmetry distinguishes access design (who may be
used) from concentration design (how value is distributed), two levers that are
easy to conflate.

%% ─────────────────────────────────────────────
\subsection{Design Implications: Governing Machine Visibility}
\label{sec:discussion:principles}

Our empirical contribution is a large, robust adoption premium attached to
agent-access admission. The conceptual step from it is a design object: the
admission boundary, in which permission architecture doubles as discovery
architecture. That coupling is partly definitional---the platform's
\texttt{allowsAgenticUsers} filter and MCP \texttt{search-actors} surface
limited-permissions, non-standby tools to machine visibility---so the
HCI questions worth asking concern what the coupling does and how to design for it. Three follow.

\emph{The boundary redistributes opportunity, not just access.} Within the same
developer, tools admitted to the agent-accessible set attract substantially more
adoption (40\% to 72\%), and two-thirds of mixed developers have a positive
ready--non-ready gap. An access rule is therefore a distributional lever: it
changes \emph{which tools get adopted}, not merely which tools are technically
admissible. Operators who treat access design as a screening question miss the
governance question of who gains and loses visibility.

\emph{Developers must design for two legibility surfaces.} A tool now has two
distinct interfaces: a visual storefront for humans and a machine-readable
surface---schema, documentation, deterministic execution---for agents. These
require different affordances, and, as API-usability research shows, each
interface must be designed explicitly for its users \cite{myers2016api}; a
developer optimizing for one audience can render the tool illegible to the
other. Dual-audience
legibility is thus not the same task done twice, but two coupled design problems
with different audiences and different success criteria---the central HCI
challenge the boundary creates.

\emph{Platforms must provide dual-audience observability.} Because
\texttt{monthly\_users} collapses human and agent traffic into one number, a
developer---or a platform---cannot tell whether a tool's adoption is
human-driven, agent-driven, or both. As machine-consumer demand grows, the
platform must expose audience-specific traffic and conversion signals, or
developers cannot make informed decisions about how to design for two audiences
whose distinct needs must be balanced---a multisided-fairness concern
\cite{burke2017multisided}.

%% ─────────────────────────────────────────────
\subsection{Limitations and Boundary Conditions}
\label{sec:discussion:limitations}

Several limitations bound the interpretation of our findings.

\emph{Single platform.} All data come from the Apify Store. Whether the
agentic-readiness adoption gap generalizes to other automation marketplaces
(Zapier, UiPath Marketplace, Bardeen) remains an open question.

\emph{Cross-sectional design and temporal mismatch.} Our core analysis uses a
single snapshot, and the treatment is measured at snapshot time while
\texttt{monthly\_users} is a trailing 30-day measure. The architectural fields
that define readiness are measured once (2026-08-27), so we cannot directly
verify that a tool held its architecture state across the full adoption window;
the store's daily whitelist flag is stable (1.4\% of actors change state), which
is suggestive but does not fully resolve the concern.

\emph{No observed agent use.} Agentic readiness describes admission to the
agent-accessible set, not realized agent use; \texttt{monthly\_users} does not
distinguish human from agent traffic, so the agent-mediated and
developer-mediated pathways cannot be separated (\autoref{sec:discussion:gap}).

\emph{Identifying subsample.} The within-developer design is identified by the
343 mixed developers---larger, older, more successful developers---so the
estimate speaks to the marketplace's established developers rather than its long
tail (\autoref{tab:mixed_developers}).

\emph{Control-group imbalance.} Only 5\% of PPE tools are non-agentic-ready, so
the comparison group is small and potentially selected; this most affects the
concentration comparison, which we address with sample-size equalization
(\autoref{sec:discussion:concentration}).

\emph{Residual within-developer selection.} Listing-time covariates barely
attenuate the within-developer coefficient; unobserved tool-level investment
remains the main alternative explanation.

%% ─────────────────────────────────────────────
\subsection{Future Work}
\label{sec:discussion:future}

Future work should study developers who selectively convert rental actors to
pay-per-event (and thus to potential agentic readiness) around the October 2026
sunset, study developers who selectively configure tools to be agentic-ready,
compare automation marketplaces (Zapier, UiPath, Bardeen), and collect
agent-side usage data to separate the agent-mediated and developer-mediated
pathways. Empirically, whether a developer optimizing for one legibility surface
renders the tool illegible to the other remains to be tested, as does how
platforms should surface and evaluate both surfaces---building on API-usability
evaluation methods \cite{myers2016api}.
